\section{Introduction to Statistics: What Are We Actually Measuring?}

Statistics helps us transform respondents' answers into information that
can be analyzed. In this study, the population consisted of Grade 12
students in Informatics II and III at SMAK Yos Sudarso Batam, while the
research sample consisted of 69 students.

The study has two main variables. The independent variable or X is the use
of gadgets as a learning medium, while the dependent variable or Y is
students' learning motivation and social interaction.

The research instrument consisted of 29 statements. A total of 15
statements were used to measure variable X, while 14 statements were used
to measure variable Y.

The questionnaire used a five-point Likert scale:

\begin{table}[htbp]
  \centering
  \singlespacing
  \caption{Likert Scale Used in the Questionnaire}
  \label{tab:likert}
  \begin{tabular}{lc}
    \toprule
    Response & Score \\
    \midrule
    Strongly Agree (SA) & 5 \\
    Agree (A) & 4 \\
    Undecided (U) & 3 \\
    Disagree (D) & 2 \\
    Strongly Disagree (SD) & 1 \\
    \bottomrule
  \end{tabular}
\end{table}

The questionnaire also contained several negatively worded statements.
During data processing, the scores for negative statements needed to be
reversed so that the direction of measurement remained consistent. The
formula used was
\[
\text{New Score} = 6 - \text{Original Score}
\]
and the procedure follows \parencite{field2018}. For example, if a
respondent gave a score of 1 to a negative statement, the score after
reverse scoring would become 5.

Therefore, statistics is not only about numbers. It begins with
understanding what is being measured, who is being studied, how the data
are collected, and how the data are transformed into information.
